\documentclass[11pt]{article}
\usepackage{amsmath,amsthm,amssymb}
\usepackage[margin=1in]{geometry}
\usepackage{enumitem}
\usepackage{booktabs}
\usepackage{hyperref}

\newtheorem{theorem}{Theorem}
\newtheorem{lemma}[theorem]{Lemma}
\newtheorem{proposition}[theorem]{Proposition}
\newtheorem{corollary}[theorem]{Corollary}
\theoremstyle{definition}
\newtheorem{definition}[theorem]{Definition}
\newtheorem{remark}[theorem]{Remark}

\title{Structural $\dim B = 5$ for $\lambda = (4, 3, 1^q)$\\
   via disjoint $\Phi$-supports and adjacent-injectivity}
\author{Clio}
\date{13 May 2026 (afternoon-3 prove session)}

\begin{document}
\maketitle

\begin{abstract}
The May-13 morning paper \texttt{2026-05-13-4-3-1n-upper-bound.tex}
proved $\dim B^{(\lambda)}_{j_0} V_\lambda \le 5$ structurally at
$q \ge 5$ for $\lambda = (4, 3, 1^q)$, with computational verification
$\dim = 5$ at $q \in \{3, 4\}$. The matching structural lower bound
$\dim \ge 5$ was deferred. The afternoon partial paper
\texttt{2026-05-13-4-3-1n-lower-bound-analysis.tex} reduced the lower
bound to two within-class projection independences and identified a
``position-of-letter contamination'' obstruction; the afternoon-2
dim2-331 paper (Remark 4.10) noted that the within-class
linear independences are individually closed by adjacent-injectivity
but concluded that this is \emph{not} sufficient to close the full
$5$-vector independence, citing a possible cross-class $c_3$-overlap
invisible to either projection. We show that this concern is
unfounded: the lower bound \emph{is} closed by adjacent-injectivity,
provided one works at the level of $V_\lambda$ directly (without
projection).

\textbf{Theorem.} For $\lambda = (4, 3, 1^q)$ with $q$ in the
structural range (the intersection of the inputs below),
\[
   \dim B^{(\lambda)}_{j_0} V_\lambda \;=\; 5.
\]

\textbf{Key idea.} The $\Phi_A^\lambda$- and $\Phi_B^\lambda$-images
have seminormal SYT-supports that are \emph{disjoint as sets of
SYTs of $\lambda$}, because the unordered cell-pairs
$\{c_1, c_3\}$ and $\{c_2, c_3\}$ differ (they share $c_3$ but have
distinct second cells). Hence
$\Phi_A^\lambda(B^{(\mu_A)}) \oplus \Phi_B^\lambda(B^{(\mu_B)})$ is a
direct sum of dimension $2 + 3 = 5$ inside $V_\lambda$. Both summands
lie in $E_{n-1}^+|_{V_\lambda}$ by construction. The outer chain
$\mathcal{O}_\lambda$ acts injectively on $E_{n-1}^+|_{V_\lambda}$ by
adjacent-injectivity (dim2-331 Theorem 3.4, applied with four
factors). Hence the five lifted vectors are linearly independent in
$V_\lambda$.

\textbf{What this resolves.} The afternoon partial paper's
Proposition 4.4 (``naive separate within-class vanishings, not
provable from position-of-$n$ alone'') was correct as a statement
about projected vectors. But the unprojected direct argument
(disjoint SYT-supports of the $\Phi_X^\lambda$-images BEFORE applying
$\mathcal{O}_\lambda$) bypasses the contamination issue entirely.
The tracer-pair Hoefsmit computation is not needed; adjacent-injectivity
suffices.

\textbf{What this generalises.} An r=2 dim-additivity statement: for
any rank-zero shape $\lambda$ with exactly two length-preserving
corners $c_1, c_2$ and column-$1$ bottom $c_3$, where the SRD-style
reduction formula gives
$B^{(\lambda)} = \mathcal{O}_\lambda[\Phi_A^\lambda(B^{(\mu_A)}) +
\Phi_B^\lambda(B^{(\mu_B)})]$, the dimensions add:
$\dim B^{(\lambda)} = \dim B^{(\mu_A)} + \dim B^{(\mu_B)}$.
The proof is shape-agnostic; the inputs are the reduction formula
and the two inner-bound dims.
\end{abstract}

\section{Setup}\label{sec:setup}

Notation as in the May-13 4-3-1n upper bound paper
(\textbf{May-13-UB}, \texttt{2026-05-13-4-3-1n-upper-bound.tex}) and
the afternoon partial paper (\textbf{LB-partial},
\texttt{2026-05-13-4-3-1n-lower-bound-analysis.tex}).
$H_q(S_n)$ is the Iwahori--Hecke algebra over $\mathbb{Q}(q)$ with
generators $T_1, \ldots, T_{n-1}$ obeying $(T_i - q)(T_i + 1) = 0$
and the standard braid relations. $V_\lambda$ is the seminormal cell
module for $\lambda \vdash n$ in asymmetric Murphy normalisation
($b' = 1$). $E_j^\pm$ denote the $q$- and $(-1)$-eigenspaces of
$T_j$. For $1 \le k \le n$,
\[
   R'_k \;:=\; (T_{k-1}{+}1)\cdots(T_1{+}1),
   \qquad
   B^{(\lambda)}_j V_\lambda \;:=\; R'_j\cdots R'_n V_\lambda.
\]

\paragraph{Throughout this paper.} $\lambda := (4, 3, 1^q)$ with
$q \ge q_*$ (with $q_*$ described below), $n := q + 7$,
$\ell := q + 2$, $j_0 := \ell + 1 = q + 3$. Removable corners
\[
   c_1 = (1, 4), \quad c_2 = (2, 3), \quad c_3 = c_* = (q+2, 1).
\]
Inner partitions (corner-pair removal):
\begin{align*}
   \mu_A &:= \lambda \setminus \{c_1, c_3\} = (3, 3, 1^{q-1}),
      &|\mu_A| = q+5, \\
   \mu_B &:= \lambda \setminus \{c_2, c_3\} = (4, 2, 1^{q-1}),
      &|\mu_B| = q+5.
\end{align*}
Outer chain:
\[
   \mathcal{O}_\lambda \;:=\;
   (T_{q+2}{+}1)(T_{q+3}{+}1)(T_{q+4}{+}1)(T_{q+5}{+}1).
\]
Four factors with indices $\{q+2, q+3, q+4, q+5\} = \{\ell, \ell+1,
n-3, n-2\}$.

\paragraph{Reduction formula (May-13-UB).}
\begin{equation}\label{eq:reduction}
   B^{(\lambda)}_{j_0} V_\lambda
   \;=\;
   \mathcal{O}_\lambda \Big[
      \Phi_A^\lambda\big(B^{(\mu_A)}_{j_0(\mu_A)} V_{\mu_A}\big)
      +
      \Phi_B^\lambda\big(B^{(\mu_B)}_{j_0(\mu_B)} V_{\mu_B}\big)
   \Big].
\end{equation}
The $\mu_C$-piece ($\Phi_C^\lambda$, removing $\{c_1, c_2\}$) and the
same-row piece ($S_1$ at row $1$, $S_2$ at row $2$) vanish structurally
for all $q \ge 3$ (May-13-UB Lemmas 3.5 and 3.6).

\paragraph{The branching injections $\Phi_X^\lambda$.}
For each corner pair $\{c_X^-, c_X^+\} \in \{\{c_1, c_3\},
\{c_2, c_3\}, \{c_1, c_2\}\}$, the map $\Phi_X^\lambda :
V_{\mu_X} \hookrightarrow V_\lambda$ sends a seminormal basis
vector $v_{T'}^{(\mu_X)}$ to the $+q$-eigenvector of the
$T_{n-1}$-$2$-block on the two SYT extensions $T_A^X, T_B^X$ of $T'$
obtained by placing $\{n-1, n\}$ at $\{c_X^-, c_X^+\}$:
\[
   \Phi_X^\lambda(v_{T'}^{(\mu_X)})
   \;=\;
   v_{T_A^X} + \tau_X v_{T_B^X},
\]
where $\tau_X$ is the Hoefsmit off-diagonal coefficient determined by
the axial distance of the pair (May-13-UB Theorem 2.2, item 1; cf.
SRD-331 Lemma 3.2).

\paragraph{Inner-bound dims.}
\begin{itemize}[leftmargin=2em,topsep=0.2em,itemsep=0.2em]
\item $\dim B^{(\mu_A)}_{j_0(\mu_A)} V_{\mu_A} = 2$.
\emph{Structural at $q - 1 \ge 4$, i.e., $q \ge 5$} (dim2-331
Theorem 1.3, applied to $\mu_A = (3, 3, 1^{q-1})$).
Computational at $q \in \{3, 4\}$ (SRD-331 Section 5).
\item $\dim B^{(\mu_B)}_{j_0(\mu_B)} V_{\mu_B} = 3$.
\emph{Structural at $|\mu_B| \ge 8$, i.e., $q \ge 3$}
(May-12-Eve3 paper \texttt{2026-05-12-evening-4-2-1n-dim3.tex},
amended with the May-13 SRD-421 same-row-dies closure,
\texttt{2026-05-13-same-row-dies-421.tex}).
\end{itemize}

\paragraph{Range of structural validity.}
The combination of inputs is structural at $q \ge q_*$, with the
binding constraint being the dim2-331 input on $\mu_A$, giving
$q_* = 5$. At $q \in \{3, 4\}$, the inner-bound dims are
computational, hence the theorem reduces to a computational statement
already verified by May-13-UB Section 5.

\section{The two key lemmas}\label{sec:lemmas}

\subsection{Disjoint $\Phi$-supports at the SYT level}

\begin{lemma}[Disjoint SYT supports]\label{lem:disjoint-supp}
For every $v \in V_{\mu_A}$ and every $w \in V_{\mu_B}$, the
seminormal supports of $\Phi_A^\lambda(v)$ and $\Phi_B^\lambda(w)$ in
$V_\lambda$ are disjoint as sets of SYTs of $\lambda$.
Consequently,
\[
   \Phi_A^\lambda(V_{\mu_A}) \;\cap\; \Phi_B^\lambda(V_{\mu_B})
   \;=\; \{0\}
\]
inside $V_\lambda$.
\end{lemma}

\begin{proof}
By the explicit description of $\Phi_X^\lambda$
(\S\ref{sec:setup}), every SYT $T$ in the seminormal support of
$\Phi_X^\lambda(v_{T'}^{(\mu_X)})$ has the two letters $\{n-1, n\}$
occupying the cell pair $\{c_X^-, c_X^+\}$ (in one of two orders).
Extending by linearity over $v \in V_{\mu_X}$, every SYT in the
seminormal support of $\Phi_X^\lambda(v)$ has $\{n-1, n\}$ at
$\{c_X^-, c_X^+\}$ as an unordered cell pair.

For $X = A$: cell pair is $\{c_1, c_3\}$.
For $X = B$: cell pair is $\{c_2, c_3\}$.
These two unordered pairs are distinct subsets of cells of $\lambda$
(they share $c_3$ but differ in the other element, $c_1$ vs $c_2$).

Hence no single SYT $T$ of $\lambda$ can lie in both supports: in $T$,
the unordered cell pair $\{T^{-1}(n-1), T^{-1}(n)\}$ is determined,
and it cannot equal both $\{c_1, c_3\}$ and $\{c_2, c_3\}$.

The seminormal basis $\{v_T : T \in \mathrm{SYT}(\lambda)\}$ is
linearly independent in $V_\lambda$. Hence any non-zero linear
combination supported on the A-set is linearly independent of any
non-zero linear combination supported on the B-set. So
$\Phi_A^\lambda(V_{\mu_A}) \cap \Phi_B^\lambda(V_{\mu_B}) = \{0\}$.
\end{proof}

\subsection{Image lies in $E_{n-1}^+$}

\begin{lemma}[$\Phi_X^\lambda$ lands in $E_{n-1}^+$]\label{lem:image-in-Eplus}
For each $X \in \{A, B\}$ and every $v \in V_{\mu_X}$,
$\Phi_X^\lambda(v) \in E_{n-1}^+|_{V_\lambda}$.
\end{lemma}

\begin{proof}
By construction (\S\ref{sec:setup}),
$\Phi_X^\lambda(v_{T'}^{(\mu_X)})$ is the $+q$-eigenvector of the
$T_{n-1}$-action on the $2$-block $\langle v_{T_A^X}, v_{T_B^X}\rangle$
of $V_\lambda$, where the two SYTs are obtained by placing $\{n-1, n\}$
at $\{c_X^-, c_X^+\}$ in the two valid orders. The axial distance
$d_X$ of this $2$-block is non-zero (the cells $c_X^-, c_X^+$ lie in
distinct rows and columns), so the $2$-block is non-degenerate and the
$+q$-eigenvector exists and is unique up to scalar. Hence
$\Phi_X^\lambda(v_{T'}^{(\mu_X)}) \in E_{n-1}^+|_{V_\lambda}$. Extending
by linearity: $\Phi_X^\lambda(V_{\mu_X}) \subseteq E_{n-1}^+|_{V_\lambda}$.
\end{proof}

\subsection{Adjacent-injectivity of the outer chain}

\begin{theorem}[Adjacent injectivity chain; dim2-331 Theorem 3.4]
\label{thm:adj-inj}
Let $V$ be any $H_q(S_n)$-module, and let
$a \le b \le n - 2$. The composition
\[
   (T_a{+}1)(T_{a+1}{+}1)\cdots(T_b{+}1)\big|_{E_{b+1}^+(V)}
   \;:\;
   E_{b+1}^+(V) \;\to\; E_a^+(V)
\]
is injective.
\end{theorem}

\begin{proof}
This is dim2-331 Theorem 3.4 (also dim2-331 Theorem 5.5). The proof is
an iterated application of the basic adjacent-injectivity proposition:
each $(T_j+1)$ maps $E_{j+1}^+(V)$ injectively into $E_j^+(V)$, by the
combination of image-landing ($(T_j+1)V \subseteq E_j^+$, via the
quadratic relation) and adjacent-disjointness
($E_j^-(V) \cap E_{j+1}^+(V) = 0$, May-19 Phase A swapped form
\texttt{2026-05-19-phase-a-general.tex} Lemma~3.3).
\end{proof}

\begin{corollary}[Outer chain is injective on $E_{n-1}^+$]\label{cor:O-inj}
The operator $\mathcal{O}_\lambda$ acts injectively on
$E_{n-1}^+|_{V_\lambda}$, with image in
$E_\ell^+|_{V_\lambda} = E_{q+2}^+|_{V_\lambda}$.
\end{corollary}

\begin{proof}
Apply Theorem~\ref{thm:adj-inj} with $a = q+2$, $b = q+5 = n-2$
(four factors).
\end{proof}

\section{The main theorem}\label{sec:main}

\begin{theorem}[Main: lower bound matches upper bound]\label{thm:main}
For $\lambda = (4, 3, 1^q)$ with $q \ge q_*$ (where $q_* = 5$ is the
binding structural threshold from the dim2-331 input on $\mu_A$),
\begin{equation}\label{eq:main}
   \dim B^{(\lambda)}_{j_0} V_\lambda \;=\; 5.
\end{equation}
\end{theorem}

\begin{proof}
\emph{Upper bound:} $\dim B^{(\lambda)} \le 5$ by May-13-UB
Theorem 1.1.

\emph{Lower bound:} Define
\[
   W \;:=\; \Phi_A^\lambda\!\big(B^{(\mu_A)} V_{\mu_A}\big)
      \;+\;
      \Phi_B^\lambda\!\big(B^{(\mu_B)} V_{\mu_B}\big)
   \;\subseteq\; V_\lambda.
\]
We compute $\dim W$ and then $\dim \mathcal{O}_\lambda W$:

\textbf{Step 1.} $W$ is a direct sum of dim $2 + 3 = 5$.
By Lemma~\ref{lem:disjoint-supp},
$\Phi_A^\lambda(V_{\mu_A}) \cap \Phi_B^\lambda(V_{\mu_B}) = \{0\}$,
and in particular
$\Phi_A^\lambda(B^{(\mu_A)}) \cap \Phi_B^\lambda(B^{(\mu_B)}) = \{0\}$.
Hence $W = \Phi_A^\lambda(B^{(\mu_A)}) \oplus \Phi_B^\lambda(B^{(\mu_B)})$
is a direct sum. Each $\Phi_X^\lambda$ is injective on $V_{\mu_X}$
(it sends distinct seminormal basis vectors to vectors with disjoint
$2$-element supports), so the inner dims lift:
\[
   \dim W \;=\; \dim B^{(\mu_A)} + \dim B^{(\mu_B)} \;=\; 2 + 3 \;=\; 5.
\]

\textbf{Step 2.} $W \subseteq E_{n-1}^+|_{V_\lambda}$.
By Lemma~\ref{lem:image-in-Eplus}, each of
$\Phi_A^\lambda(B^{(\mu_A)})$ and $\Phi_B^\lambda(B^{(\mu_B)})$ lies in
$E_{n-1}^+|_{V_\lambda}$. Their sum (direct or not) also lies there.

\textbf{Step 3.} $\mathcal{O}_\lambda$ is injective on $W$.
By Corollary~\ref{cor:O-inj}, $\mathcal{O}_\lambda$ is injective on
all of $E_{n-1}^+|_{V_\lambda}$, hence on the subspace $W$.

\textbf{Step 4.} Conclude. By Steps 1, 3:
\[
   \dim \mathcal{O}_\lambda W \;=\; \dim W \;=\; 5.
\]
By the reduction formula~\eqref{eq:reduction},
$B^{(\lambda)}_{j_0} V_\lambda = \mathcal{O}_\lambda W$, so
$\dim B^{(\lambda)} \ge 5$.

Combined with the upper bound: $\dim B^{(\lambda)} = 5$.
\end{proof}

\section{Why the afternoon partial paper got stuck}\label{sec:partial}

The afternoon partial paper proceeded via the position-of-$n$
projection $\pi^{(c)}_n$ (for $c \in \{c_1, c_2, c_3\}$). Its
Theorem 3.2 reduced the full $5$-vector independence to two
within-class projection independences:
\begin{enumerate}[leftmargin=2em,topsep=0.2em,itemsep=0.1em,label=(\alph*)]
\item $\{\pi^{(c_1)}_n F_A^1, \pi^{(c_1)}_n F_A^2\}$ linearly
independent.
\item $\{\pi^{(c_2)}_n F_B^1, \pi^{(c_2)}_n F_B^2, \pi^{(c_2)}_n F_B^3\}$
linearly independent.
\end{enumerate}
Both projections are taken AFTER applying $\mathcal{O}_\lambda$. The
issue, as flagged in dim2-331 Remark 4.10: the projected vectors are
not in $E_{n-1}^+$ (the seminormal projection
$\pi^{(c)}_n$ breaks the $T_{n-1}$-2-block $+q$-eigenvector structure),
so adjacent-injectivity does not apply to the projected vectors. The
``cross-class $c_3$-overlap'' concern: a linear combination
$\sum \alpha_i F_A^i + \sum \beta_j F_B^j = 0$ could carry weight on
SYTs with $n$ at $c_3$ (where both $\Phi_A$- and $\Phi_B$-pieces
project non-trivially), and this overlap could be invisible to either
$\pi^{(c_1)}_n$ or $\pi^{(c_2)}_n$ alone.

\paragraph{Resolution.} The concern is real for the projected approach
but does \emph{not} obstruct the direct approach. Working at the
level of $V_\lambda$ (no projection):

\begin{itemize}[leftmargin=2em,topsep=0.2em,itemsep=0.2em]
\item Before applying $\mathcal{O}_\lambda$: the lifts
$\Phi_A^\lambda(v)$ and $\Phi_B^\lambda(w)$ have seminormal supports
on SYTs distinguished by the position of \emph{both} $n-1$ AND $n$,
not just $n$. The cell-pair $\{n-1, n\}$ is at $\{c_1, c_3\}$ for
A-lifts and at $\{c_2, c_3\}$ for B-lifts; these are distinct
unordered pairs, hence disjoint SYT supports
(Lemma~\ref{lem:disjoint-supp}).
\item After applying $\mathcal{O}_\lambda$: the position of $n-1$ is
moved by the factors $T_{n-2}, T_{n-3}$ in $\mathcal{O}_\lambda$, so
the SYT supports of $F_A^i$ and $F_B^j$ can overlap (in particular at
SYTs with $n$ at $c_3$). \emph{However}, this overlap is the image of
two linearly independent subspaces $\Phi_A^\lambda(B^{(\mu_A)})$ and
$\Phi_B^\lambda(B^{(\mu_B)})$ under the injective map
$\mathcal{O}_\lambda|_{E_{n-1}^+}$, so the images remain linearly
independent.
\end{itemize}

\paragraph{Moral.} The decoupling argument should be performed
\emph{before} the outer chain acts. Position-of-$(n-1)$ at the
\emph{inner} level (i.e., on the $\Phi_X^\lambda$-supports inside
$V_\lambda$) is automatic and clean. The outer chain then acts on
the already-decoupled direct sum, and injectivity preserves the
decoupling.

\paragraph{Comparison to position-of-letter contamination.} The
afternoon partial paper's ``position-of-letter contamination at
recursion depth 3'' is a real obstruction to lifting position-of-$(n-2)$
\emph{within} a single class through three layers of recursion. That
contamination prevents the within-A separation
$\{F_A^1, F_A^2\}$ from being detected by any single position-of-letter
projection at the $\lambda$ level. But within-class separation is
\emph{also} closed by adjacent-injectivity applied to
$\Phi_A^\lambda(B^{(\mu_A)})$ alone (dim2-331 Remark 4.10's first
sentence). The afternoon-2 author noted this within-class fact but
incorrectly concluded that it didn't combine with A-vs-B decoupling
to close the full lower bound. The present paper observes that the
A-vs-B decoupling is also available at the inner level (disjoint
$\Phi$-supports), and both pieces of decoupling chain together
through the injective outer chain.

\section{Computational verification}\label{sec:comp}

Script:
\texttt{\textasciitilde/projects/scratch/2026-05-13-431-lower-closed/verify\_431\_lower.py}.
Mod-$P$ arithmetic at $P = 100\,003$, $q = 11$, for
$\lambda = (4, 3, 1^q)$ at $q \in \{3, 4\}$:

\begin{itemize}[leftmargin=2em,topsep=0.3em,itemsep=0.2em]
\item \textbf{Disjoint SYT supports.} For each $q$, the number of SYTs
of $\lambda$ with $\{n-1, n\}$ at $\{c_1, c_3\}$ and at
$\{c_2, c_3\}$:
\begin{center}
\begin{tabular}{rrrr}
\toprule
$q$ & $n$ & A-support & B-support \\
\midrule
$3$ & $10$ & $112$ & $180$ \\
$4$ & $11$ & $240$ & $378$ \\
\bottomrule
\end{tabular}
\end{center}
Intersection is empty in both cases (consistent with
Lemma~\ref{lem:disjoint-supp}).

\item \textbf{Adjacent disjointness.} $E_j^- \cap E_{j+1}^+ = 0$
verified at $j \in \{q+2, q+3, q+4, q+5\}$ for both
$q \in \{3, 4\}$.

\item \textbf{Outer chain injectivity.}
$\mathrm{rank}(\mathcal{O}_\lambda \cdot E_{n-1}^+) = \dim E_{n-1}^+$
in both cases ($245$ at $q=3$, $470$ at $q=4$).

\item \textbf{Direct dim.} $\dim B^{(\lambda)}_{j_0} V_\lambda = 5$ at
each $q \in \{3, 4\}$, matching the conjectured value (May-13-UB
Section~5).
\end{itemize}

\section{The r=2 dim-additivity meta-theorem}\label{sec:r2}

The proof of Theorem~\ref{thm:main} uses only three shape-sensitive
inputs and one universal fact:

\paragraph{Shape inputs.}
\begin{itemize}[leftmargin=2em,topsep=0.2em,itemsep=0.1em]
\item The SRD-style reduction formula
$B^{(\lambda)} = \mathcal{O}_\lambda[\Phi_A^\lambda(B^{(\mu_A)})
+ \Phi_B^\lambda(B^{(\mu_B)})]$, which holds when same-row-dies
applies to the same-row pieces of $\lambda$ and the $\Phi_C^\lambda$
piece vanishes ($j$-formula leakage on the $(c_1, c_2)$ block).
\item Non-degeneracy of the $T_{n-1}$-$2$-blocks for both corner pairs
$\{c_X^-, c_X^+\}$, $X \in \{A, B\}$ (automatic when corners are at
distinct row+column).
\item The inner-bound dims $\dim B^{(\mu_A)}$, $\dim B^{(\mu_B)}$.
\end{itemize}

\paragraph{Universal fact.} Adjacent-injectivity
(Theorem~\ref{thm:adj-inj}), applied to the four-factor outer chain.

\paragraph{Combinatorial fact.}
Disjointness of the unordered cell-pairs $\{c_A^-, c_A^+\}$ and
$\{c_B^-, c_B^+\}$ (automatic when $c_A^- \ne c_B^-$).

\begin{theorem}[r=2 dim-additivity]\label{thm:r2-additivity}
Let $\lambda \vdash n$ be a rank-zero partition with exactly two
length-preserving corners $c_1, c_2$ and column-$1$ bottom $c_*$,
satisfying:
\begin{enumerate}[leftmargin=2em,topsep=0.2em,itemsep=0.1em]
\item The SRD-style reduction formula~\eqref{eq:reduction} holds
for $\lambda$ with $\mathcal{O}_\lambda = (T_\ell{+}1)\cdots(T_{n-2}{+}1)$.
\item Both corner-pair $T_{n-1}$-$2$-blocks $\{c_1, c_*\}$ and
$\{c_2, c_*\}$ are non-degenerate ($|d| \ge 2$ axial distance).
\end{enumerate}
Then
\[
   \dim B^{(\lambda)}_{j_0} V_\lambda
   \;=\;
   \dim B^{(\mu_A)}_{j_0(\mu_A)} V_{\mu_A}
   \;+\;
   \dim B^{(\mu_B)}_{j_0(\mu_B)} V_{\mu_B},
\]
where $\mu_A := \lambda \setminus \{c_1, c_*\}$,
$\mu_B := \lambda \setminus \{c_2, c_*\}$.
\end{theorem}

\begin{proof}
Repeat the proof of Theorem~\ref{thm:main} verbatim. The disjoint
SYT-supports lemma (Lemma~\ref{lem:disjoint-supp}) generalises to any
two corner pairs $\{c_X^-, c_X^+\}$ that differ as unordered sets;
the $\{c_1, c_*\}$ vs $\{c_2, c_*\}$ pairs differ in their second
element (since $c_1 \ne c_2$). The image-in-$E_{n-1}^+$ lemma
(Lemma~\ref{lem:image-in-Eplus}) uses only non-degeneracy of the
corner-pair blocks (hypothesis (ii)). The adjacent-injectivity
chain (Corollary~\ref{cor:O-inj}) is universal and applies to any
outer chain.
\end{proof}

\begin{remark}[Generalisation to higher $r$]\label{rem:higher-r}
The r=2 argument generalises immediately to r=3 and beyond, with one
caveat. For $r$ length-preserving corners $c_1, \ldots, c_r$ plus the
column-$1$ bottom $c_*$, the reduction formula has $r$ corner-pair
$\Phi_X$-pieces (each removing $\{c_i, c_*\}$) plus
$\binom{r}{2}$ pure-length-pres $\Phi_{ij}$-pieces (each removing
$\{c_i, c_j\}$). The unordered cell pairs are:
\[
   \{\{c_1, c_*\}, \ldots, \{c_r, c_*\}\}
   \;\cup\;
   \{\{c_i, c_j\} : 1 \le i < j \le r\}.
\]
These $r + \binom{r}{2}$ pairs are pairwise distinct (no two share
both cells), so the seminormal supports of the corresponding
$\Phi$-images are pairwise disjoint. If the SRD-style reduction
restricts the contributing pieces to a subset $\mathcal{S}$ (with the
others vanishing by $j$-formula leakage or same-row-dies), then
\[
   \dim B^{(\lambda)} \;=\; \sum_{X \in \mathcal{S}}
   \dim B^{(\mu_X)}.
\]
This formalises and proves a long-conjectured ``additivity'' formula
for $\dim B$ in the stable regime.
\end{remark}

\section{Honest status}\label{sec:honest}

\paragraph{What we proved structurally.}
$\dim B^{(\lambda)}_{j_0} V_\lambda = 5$ for $\lambda = (4, 3, 1^q)$
at $q \ge 5$ (structural range of the dim2-331 input on $\mu_A$).

\paragraph{What we corrected.}
The dim2-331 Remark~4.10 and the afternoon partial paper's
``tracer-pair-required'' conclusion. Adjacent-injectivity \emph{is}
sufficient to close the $5$-vector independence, provided one
applies it to the unprojected direct sum
$\Phi_A^\lambda(B^{(\mu_A)}) \oplus \Phi_B^\lambda(B^{(\mu_B)})$ at
the $V_\lambda$ level. The afternoon partial paper's position-of-letter
contamination obstruction is real but is an obstruction only to the
projected approach.

\paragraph{What this generalises.}
An r=2 dim-additivity theorem (Theorem~\ref{thm:r2-additivity}),
parallel to the dim2-331 r=1 dim-equality theorem. Together, the
two theorems give a complete recursive computation of $\dim B$ for
any rank-zero shape with SRD-style reduction, in the stable regime.

\paragraph{What remains.}
\begin{itemize}[leftmargin=2em,topsep=0.2em,itemsep=0.1em]
\item Confirming the structural range of the SRD-421 input on
$\mu_B$: the May-12-Eve3 + May-13 SRD-421 combination is structural at
$|\mu_B| \ge 8$, but the precise lower threshold of $q$ for the
$(4, 3, 1^q)$ application should be cross-checked against the actual
chain of inputs.
\item Verifying Remark~\ref{rem:higher-r} explicitly on a multi-corner
r=3 shape (e.g.\ $(5, 3, 2, 1^*)$ or $(5, 4, 1^*)$). The closed-form
$\dim B = f^{\lambda^{(\ge 2)}}$ from the May-13 SYT-formula paper
should now be derivable from r=2 (or r=3) additivity plus the inner
chain reduction.
\end{itemize}

\paragraph{Connection to the closed-form SYT theorem.}
The May-13 \texttt{2026-05-13-closed-form-syt.tex} paper conjectured
$\dim B^{(\lambda)} = f^{\lambda^{(\ge 2)}}$ in the stable regime.
Theorem~\ref{thm:r2-additivity} (and its r=3+ generalisation in
Remark~\ref{rem:higher-r}) suggests an inductive proof:
$f^{\lambda^{(\ge 2)}}$ satisfies the same recursion
$\dim B^{(\lambda)} = \sum_i \dim B^{(\mu_i)}$ by SYT branching, and
the r=2 additivity matches it at the base. A full structural proof of
the closed-form SYT theorem may now be within reach, provided the
r=$r$ additivity holds in its claimed generality for all shapes in
the stable regime.

\end{document}
